% M10, 11.09.2026: Leserführung, Kapiteltrennung und belegte Reichweite.
% Architekturrolle und drei Zustandsarten; E47-01 bleibt in 4.9/Anhang A.
% Mechanik 6.4--6.6; Nachweise 7.6. Kein neuer Versuchsfall.

\section{Durable Ausführung und agentische Bedienung}
\label{sec:durable-ausfuehrung-agentische-bedienung}

Die sechste Architektursäule verbindet den unterbrechbaren Betrieb
(A9) mit der agentischen Bedienung (A10). Ein Workflow bestimmt die
Verarbeitungsschritte und ihre Übergänge. Der Steward ergänzt diese
Steuerung um die Interpretation von Nutzeranliegen und die Auswahl
passender Systemfähigkeiten. Damit bleiben die fachlichen Abläufe
explizit festgelegt, während der Zugang zu ihnen dialogisch gestaltet
ist.

Für dieses Zusammenspiel sind drei Zustandsarten auseinanderzuhalten.
Der \emph{Core} enthält die maßgeblichen Projektinformationen. Der
\emph{Ausführungszustand} hält den Fortschritt eines einzelnen Laufs
einschließlich ausstehender Antworten fest. Der
\emph{Sitzungszustand des Stewards} bewahrt den Gesprächszusammenhang.
Eine Wiederaufnahme des Laufs stellt dessen Fortsetzung her; ein
gespeicherter Dialog ermöglicht die Fortführung des Gesprächs. Keiner
dieser Vorgänge setzt für sich genommen den Core zurück oder autorisiert
neue fachliche Änderungen.

Wartet ein Lauf auf eine menschliche Entscheidung, kann sein
Ausführungszustand an einem Sicherungspunkt gespeichert werden. Nach
Einreichung der Antwort setzt er von dort fort. Gegen unbeabsichtigte
Wiederholung bereits vollzogener Wirkungen sind zusätzlich
projektspezifische Prüfungen erforderlich. Die Architektur sieht solche
Prüfungen für ausgewählte Anwendungspfade vor; sie beansprucht keine
allgemeine Garantie genau einmaliger externer Ausführung. Die Mechanik
und ihre Betriebsgrenzen behandelt
Abschnitt~\ref{sec:realisierung-durabilitaet}.

Der Steward kann Projektinformationen und Laufstatus abfragen,
Verarbeitung anstoßen und anstehende Entscheidungen im Dialog
vermitteln. So kann ein fachlicher Präzisierungswunsch zum Eingang
eines regulären Fortschreibungslaufs werden. Die Bedienschicht erhält
dafür einen begrenzten Werkzeugraum; sie besitzt keinen unmittelbaren
Schreibzugriff auf den Core. Ihre fachlichen Wirkungen laufen über
die vorgesehenen Prüf-, Entscheidungs- und Speicheroperationen.
Werkzeugzustimmung und fachliche Autorisierung bleiben dabei getrennt
(Abschnitt~\ref{sec:governance-mutation}): Die Erlaubnis, einen Ablauf
zu starten, ist noch keine Freigabe seines späteren Änderungsvorschlags.

Diese Rollenverteilung verbindet agentische Handlungsauswahl mit
kontrollierter Workflowausführung. Ihre technische Umsetzung zeigt
Abschnitt~\ref{sec:realisierung-steward}; die Wiederaufnahmebelege und
zwei dokumentierte Steward-Aufgaben werden in
Abschnitt~\ref{sec:eval-resume} eingeordnet. Daraus folgt keine
Bewertung der Bedienqualität. Der folgende Abschnitt führt die
bisherigen Architekturbausteine an den beiden End-to-End-Pfaden
zusammen.
